% !TeX spellcheck = ru_RU
% !TEX root = vkr.tex

\section{Постановка задачи}
\label{sec:task}

Цель выпускной квалификационной работы — разработка учебного программного комплекса, объединяющего виртуальный RISC‑V процессор, минимальную UNIX‑подобную операционную систему и легковесный компилятор для языка С. Комплекс предназначен для наглядной демонстрации жизненного цикла разработки и функционирования программ - от компиляции исходного кода до выполнения на процессоре. Комплекс предполагается использовать в рамках университетских курсов по архитектуре ЭВМ и системному программированию.

Для осуществления этой цели были сформулированы следующие задачи.

\begin{enumerate}
    \item Провести обзор открытых процессорных архитектур , легковесных UNIX‑подобных ОС и компактных С-компиляторов с точки зрения применимости в образовательных проектах; на основе обзора выработать требования к учебному стенду и обосновать выбор технологического стека.
    \item Спроектировать  архитектуру программного комплекса, определяя взаимодействие виртуального процессора с периферией, ОС и легковесным С-компилятором спроектировать средства наблюдения и отладки.
    \item Реализовать на SystemVerilog синтезируемое ядро RISC‑V RV32I с полной поддержкой целочисленной арифметики и рядом других возможностей.
    \item Выполнить адаптацию и сборку исходного кода UNIX‑подобной ОС для созданного процессорного ядра, обеспечив корректную загрузку и выполнение ядра на виртуальном процессоре.
    \item Адаптировать легковесный компилятор языка С для работы внутри гостевой ОС, обеспечив возможность компиляции и выполнения пользовательских C-программ на моделируемом RISC-V процессоре.
    \item Провести тестирование и верификацию всех компонентов комплекса на задачах управления процессами, памятью и вводом‑выводом.
    \item Подготовить учебно‑методические материалы для использования комплекса в учебном процессе.
\end{enumerate}
